\subsection{一致空间中的紧集}

下面这个关于任意一致空间的命题，是第一章 § 9 第 2 目命题 2 对紧空间的更精确形式。

命题 4 在一致空间 X 中，设 A 是紧集，B 是闭集，且 A∩B=∅。则存在 X 的一致邻域 V，使 V(A) 与 V(B) 不相交。

若命题不真，则当 V 取遍 X 的对称一致邻域时，诸集合 \(A \cap \mathring{V}(B)\) 都不空；于是这些集合将在 A 上构成一个滤子基，它有一个聚点 \(x_0 \in A\)。因此对每个对称一致邻域 \(V\)（\(X\) 的），\(\dot{V}(x_0)\) 将与 \(B\) 相交，从而由于 \(B\) 是闭的，应有 \(x_0 \in B\)，与假设矛盾。

推论 设 A 是一致空间 X 中的紧集；则当 V 取遍 X 的一致邻域时，诸集合 V(A) 构成 A 的一个基本邻域系。

设 U 是 A 的任一开邻域，则 B = CU 是闭的且不与 A 相交；于是由命题 4，存在一致邻域 V 使 \(\mathrm{V}(\mathrm{A}) \cap \mathrm{V}(\mathrm{B}) = \varnothing\)，因此 \(\mathrm{V}(\mathrm{A}) \subset \mathrm{U}\)。

